\scanpage{59}
\mspage{29}
prets as the group \inserted{\unsure{compatible} with} of $\unsure{U|\deleted{K}|\deleted{U}}$-automorphisms of the scheme $\widetilde U$, and $E_U\longrightarrow E_K$ as the obvious homomorphism of passage to the quotient (\uline{NB} $\bar K$ is identified with the algebraic closure of $K$ in $H^0(\widetilde U,\underline{\mathcal O}_{\widetilde U})$, hence $E_U$ acts on $\operatorname{Spec}\bar K$ in a natural way\ldots) A section of $E_U$ over $E_K$ is therefore an action of $E_K$ on $\widetilde U$, compatible with its action on $\widetilde U$, and continuous in a suitable sense (probably the pointwise-discr\unsure{e|e}te convergence on the \deleted{$\unsure{\bar U_i^{a}|\widetilde U_i^{a}|\bar U_i^{\alpha}|\widetilde U_i}$} $\bar U_i$ finite over $\bar U$ between $\widetilde U$ and $\bar U$\ldots).

Let us therefore consider the

\uline{Conjecture} \uline{B} Every such action of $\Gamma$ on $\widetilde U$ admits \uline{in $X'$}$=\widehat{\widetilde U}$ one fixed point and only one.
\margin{\figure{\begin{tikzpicture}\draw (0,0)--(0,-3.1);\end{tikzpicture}}}
\margin{and \unsure{even|even} the induced action of $\Gamma^\circ$ must have one fixed point and only one --- cf below.}

This therefore means

a) If there is a fixed point at finite distance ie. $\widetilde x\in\widetilde U^\Gamma$, then 1°) The image of $\widetilde x$ in $U$ is uniquely determined --- this is essentially theorem 3 of $\S1$ (two distinct points of $U(K)$ define \unsure{distinct conjugacy classes of splittings|conjugacy classes of distinct splittings}) and
\margin{58}
